\section{Limitations}\label{sec:limitations}

We state these as scope decisions and known weaknesses, in roughly decreasing order of importance for interpreting the results.

\begin{enumerate}
\item \textbf{Small sample, very small at-risk subsets.} There are 50 records; false reassurance can only arise on records whose truth is abnormal, and there are only 9 (Sweep~1) or 10 (Sweep~2) such records for heart rate and 19 usable ones for QTc. The counts show that the failure exists and where it concentrates; they are not population rates, and their confidence intervals are wide (for example 4 of 9: 95\,\% CI 19--73\,\%). The reduction from about 500 records was forced by GPU quota and time, and was decided before any result was seen.
\item \textbf{Agreement, not accuracy, for intervals.} PR, QRS, QT and QTc were obtained by applying the \emph{same} algorithm to the ground-truth signal and the digitised signal. They therefore measure agreement between two applications of one delineator, not agreement with clinician-annotated intervals, and the ground-truth QTc classification is algorithmic. The 20\,ms floor mixes digitiser error, delineator sensitivity and ground-truth sampling rate; we separated only the last.
\item \textbf{One digitiser, one renderer, one lead, simulated damage.} All results describe Open-ECG-Digitizer applied to ECG-Image-Kit pages degraded by post-processing. Real photographs combine degradations, add glare, creases, curvature and partial occlusion, and are captured by cameras with their own processing. The findings may not transfer to other digitisers or to photographs, and we make no claim that they do. Only lead~II was analysed.
\item \textbf{The digitiser was not version-pinned.} Sweeps 1 and 2 may have used different builds, and the exact versions were not recorded (Section~\ref{sec:drift}). Which records count as unreliable is therefore a property of a build, and Sweep~1's exclusions and Sweep~2's exclusions are not interchangeable. A per-test comparison showed that the sweeps agree on the 43 records readable in both (Section~
ef{sec:drift}), but this repeat shares the same 50 records, images and code, so it is not an independent replication.
\item \textbf{The QTc outcome analysis is exploratory.} It required two corrections found during analysis (excluding records already wrong on their clean image, and a 10\,ms margin), both chosen after seeing the data, and it rests on six events.
\item \textbf{Single-axis degradation.} One axis was varied at a time; real photographs vary several at once and combinations may compound.
\item \textbf{Selection by reliability.} Excluding unreliable records is a data-dependent choice; a record that fails on one build and not on another shows the choice is not neutral.
\item \textbf{Multiple comparisons and covariates.} Holm correction was applied within each metric but not across metrics, and no model adjusted for class, sex or age. Effect sizes are reported beside every significance statement for that reason.
\item \textbf{Signal fidelity was not measured.} The planned SNR tier was computed incorrectly and is not reported.
\item \textbf{Scope of the wider project.} Explanation to the patient, the safety boundary, and the application were not built or evaluated. Nothing in this paper is evidence about them or about clinical benefit.
\item \textbf{No clinical validation.} Thresholds follow published guidance for classification only. This is not a medical device and no diagnostic claim is made.
\end{enumerate}

\section{Conclusions}
In a controlled simulation of photographed paper ECGs, digitisation error was small for most degradations and large for two: severe blur and low resolution. Only these produced falsely reassuring heart rates, and they did so only at their most severe level, a result that a repeat sweep reproduced exactly (a computational repeat, not independent evidence) even though seven records changed readability. Tilt produced unreadable pictures, not wrong ones; JPEG compression, faded grids and uneven lighting did nothing measurable. Interval measurements confirmed the same pattern on top of a much larger error floor, part of which came from our own low-rate ground truth, and a QTc analogue of false reassurance was weak and not concentrated. Two lessons extend beyond this study: measure a new metric's clean-image baseline before interpreting degradation effects, and pin and record tool versions. The next steps are a version-pinned re-run on more records, a more precise interval delineation with explicit uncertainty, real photographs, a second digitiser, and only then the explanation and safeguard stages that motivated the work.

\subsection*{Work started after this manuscript}
Since this study was written up, a first step on the image-quality proxy named above as future work has begun, and is recorded here for completeness even though it was not part of the design of the study reported above and was not subject to its pre-registered scope decisions. Five signals computable from a photograph alone (no ground truth, no digitiser internals) were implemented: Laplacian-variance blur, a Hough-line-based skew estimate, an FFT high-frequency-energy resolution estimate, global intensity contrast, and a spectral JPEG-blockiness measure. On a small spike test (six demonstration pages, not the 50-record study sample, degraded with the same code used to build this study's images), every signal moved monotonically with its targeted degradation axis (Spearman $|\rho|$ 0.88--0.98 against level). Two needed correction before they worked: the skew estimate initially locked onto the ECG trace's own curved edges rather than the page's straight structural lines, fixed by requiring candidate lines to span at least 30\,\% of the image width; and the blockiness measure's first, spatial version was unusable noise, replaced by a spectral one that is monotonic but moves opposite to the naive expectation (weaker signal, not stronger, as JPEG damage increases), for a reason not yet established.

The five signals were then computed for all 1{,}550 images of the present study (not the six-page spike sample) and compared against the results already reported above, rather than left as an untested proposal. Four of the five correlated with the real measured QRS-duration error (Spearman, matched on 1{,}509 of 1{,}550 tests): blur ($\rho=-0.306$, $p=3.4\times10^{-34}$), resolution ($\rho=-0.285$, $p=1.2\times10^{-29}$), contrast ($\rho=-0.260$, $p=1.2\times10^{-24}$) and blockiness ($\rho=+0.085$, $p=9.1\times10^{-4}$); skew did not ($\rho=-0.020$, $p=0.47$). The two strongest correlations, blur and resolution, are the same two axes this study found to matter for both heart rate (Section~\ref{sec:hr1}) and intervals (Section~\ref{sec:intervalerr}), obtained here independently from measured error rather than assumed from the design. Using the same false-reassurance definition, the same 49 reliable records and the same 19 usable at-risk records (589 individual tests) reported in Section~\ref{sec:qtcoutcome}, a logistic-regression classifier fit on all five signals reached an AUC of 0.785 against the same six real events reported there -- the independent recomputation reproduced that count exactly, which is itself a check on both analyses. That AUC of 0.785 was, however, fit and evaluated on the same 589 tests it was trained on, with no held-out split, and rests on only six events, so it is optimistic and is not reported as a held-out result. A stricter follow-up held out whole recordings (49 recordings, 1{,}519 tests, repeated five-fold, cluster-bootstrap interval) and widened the target to any reading that was unreadable or wrong by more than a fixed tolerance. For heart rate off by more than 5\,bpm (99 bad readings) the five signals reached an AUC of 0.607 (95\% interval 0.49--0.82), against 0.655 for an oracle told the exact damage type and level; for larger errors the result was stronger (heart rate $>$10\,bpm 0.825, QTc $>$60\,ms 0.767, QRS $>$100\,ms 0.736). Training without one damage type and testing on it gave AUCs of 0.32--0.37 for brightness, JPEG and rotation, so the detector does not yet generalise to damage it was not built around. It is therefore reported as a promising prototype, not a working detector.
